\documentclass[10pt,journal,onecolumn]{IEEEtran}

% ---------- Packages ----------
\usepackage[T1]{fontenc}
\usepackage{newtxtext}
\usepackage{amsmath}
\usepackage{amssymb}
\usepackage{newtxmath}
\usepackage{booktabs}
\usepackage{array}
\usepackage{tabularx}
\usepackage{multirow}
\usepackage{listings}
\usepackage{xcolor}
\usepackage[hidelinks]{hyperref}

% ---------- Paragraph breaking ----------
\emergencystretch=3em
\tolerance=2000

% ---------- Code style ----------
\definecolor{codegray}{gray}{0.94}
\definecolor{codebar}{gray}{0.65}
\lstdefinestyle{srs}{
  basicstyle=\ttfamily\footnotesize,
  backgroundcolor=\color{codegray},
  breaklines=true,
  showstringspaces=false,
  frame=single,
  rulecolor=\color{codebar},
  framerule=0.3pt,
  columns=fullflexible,
  keepspaces=true,
  aboveskip=4pt,
  belowskip=4pt
}
\lstset{style=srs}

% ---------- Commands ----------
\newcommand{\reqid}[1]{\textbf{#1}}

% ---------- Title / Authors ----------
\title{Software Requirements Specification:\\[0.4em] NeuroShell Intent Recognition Engine (IRE)}
\author{\IEEEauthorblockN{NeuroShell Core Team}
\IEEEauthorblockA{Cyber Operations Research Group\\
MSc Cybersecurity Research Project R26-CS-018}}

\begin{document}

\maketitle

\begin{abstract}
This Software Requirements Specification (SRS) defines the complete functional and non-functional requirements for the NeuroShell Intent Recognition Engine (IRE), designated Component 01 of the NeuroShell Framework. The IRE is a domain-specific Natural Language Understanding (NLU) gateway that converts natural-language offensive-security operator commands into validated, schema-constrained JSON intent contracts. This document specifies the seven-stage pipeline architecture with its eleven middleware components, the hallucination and adversarial detection mechanisms, role-based access control, external interfaces, non-functional requirements, data requirements, and the empirical verification results obtained from a 900-record evaluation campaign.
\end{abstract}

\begin{IEEEkeywords}
Software Requirements Specification, Intent Recognition, Natural Language Understanding, Large Language Models, FastAPI, Cybersecurity Automation, Penetration Testing.
\end{IEEEkeywords}

% =================== Document Control ===================
\section*{Document Control}
\begin{center}
\begin{tabularx}{0.9\columnwidth}{>{\bfseries}l >{\raggedright\arraybackslash}X}
Document ID & NS-IRE-SRS-001\\
Version & 1.0\\
Date & 2026-09-01\\
Project & NeuroShell Framework --- Component 01\\
Reference & R26-CS-018 (MSc Cybersecurity Research)\\
Classification & Internal / Academic Research\\
Status & Current\\
\end{tabularx}
\end{center}

\section*{Revision History}
\begin{center}
\begin{tabularx}{0.9\columnwidth}{l l l X}
\toprule
\textbf{Version} & \textbf{Date} & \textbf{Author} & \textbf{Changes}\\
\midrule
1.0 & 2026-09-01 & NeuroShell IRE & Initial comprehensive SRS\\
\bottomrule
\end{tabularx}
\end{center}

\tableofcontents

\newpage

% =================== 1. INTRODUCTION ===================
\section{Introduction}

\subsection{Purpose}
This Software Requirements Specification (SRS) defines the complete functional and non-functional requirements for the \textbf{NeuroShell Intent Recognition Engine (IRE)}, designated Component 01 of the NeuroShell Framework. The IRE is a domain-specific Natural Language Understanding (NLU) gateway that converts natural-language offensive-security operator commands into validated, schema-constrained JSON intent contracts.

This document is intended for:
\begin{itemize}
\item \textbf{Research evaluators} assessing the IRE's design and empirical performance;
\item \textbf{Developers} maintaining or extending the pipeline;
\item \textbf{Downstream consumers} (e.g., Component 02 Dynamic Planner) integrating with the IRE API;
\item \textbf{Examiners} reviewing the MSc Cybersecurity research project (ref.~R26-CS-018).
\end{itemize}

\subsection{Scope}

\subsubsection{Product Scope}
The IRE is a \textbf{self-contained FastAPI microservice} that:
\begin{enumerate}
\item Accepts natural-language offensive-security commands via a REST API;
\item Applies multi-layer validation (adversarial detection, schema, regex, network architecture, scope);
\item Produces a validated \texttt{IntentSchema} or \texttt{PlannerContract} JSON output;
\item Enforces role-based access control, session context, and hallucination detection;
\item Maintains a durable, privacy-preserving audit log.
\end{enumerate}

\subsubsection{In Scope}
\begin{itemize}
\item All source code under \texttt{src/} (pipeline, middleware, validation, inference, schemas, API, audit, utilities);
\item Configuration under \texttt{config/} (settings, features, RBAC policy, enforcement policy, LoRA config);
\item Grounding data under \texttt{data/} (CVE grounding, alias map);
\item Evaluation harness under \texttt{eval/} (diagnose, metrics, experiment scripts);
\item Test suite under \texttt{tests/} (unit and integration);
\item Deployment artifacts (Dockerfile, docker-compose.yml, scripts).
\end{itemize}

\subsubsection{Out of Scope}
\begin{itemize}
\item Component 02 (Dynamic Planner), Component 03 (AEERE Executor), and Component 04 (AVAE Vulnerability Analyzer) --- these are downstream consumers referenced only for interface-design justification;
\item Frontend UI (separate project);
\item Network infrastructure and physical security;
\item LoRA fine-tuning execution (Phase 2, planned but not yet implemented).
\end{itemize}

\subsection{Definitions, Acronyms, Abbreviations}
\begin{table}[h]
\caption{Definitions, Acronyms, and Abbreviations.}
\label{tab:definitions}
\begin{tabularx}{\columnwidth}{>{\ttfamily}l X}
\toprule
\textbf{Term} & \textbf{Definition}\\
\midrule
IRE & Intent Recognition Engine\\
NLU & Natural Language Understanding\\
LLM & Large Language Model\\
CVE & Common Vulnerabilities and Exposures\\
RBAC & Role-Based Access Control\\
CIDR & Classless Inter-Domain Routing\\
RFC-1918 & Address Allocation for Private Internets\\
SRS & Software Requirements Specification\\
Tier-1 & Single-sample LLM inference path\\
Tier-2 & Conditional N-sample self-consistency path\\
DAcc & Decision Accuracy\\
CFAcc & Contract-Field Accuracy\\
FP & False Positive (incorrectly blocked well-formed input)\\
S5/S6/S6B/S7 & Pipeline validation stages (Schema/Regex/Network/Scope)\\
M1--M11 & Pipeline middleware components\\
PlannerContract & Machine-readable JSON hand-off object for Component 02\\
Grounding Hash & SHA-256 digest of CVE and alias datasets for cache invalidation\\
\bottomrule
\end{tabularx}
\end{table}

\subsection{References}
\begin{thebibliography}{10}
\bibitem{R1} NeuroShell IRE README. \emph{README.md}, NeuroShell IRE Repository.
\bibitem{R2} IRE Research Extraction. \emph{IRE-Doc/research-extraction.md}.
\bibitem{R3} Experiment Results (900-record). \emph{IRE-Doc/NeuroShell\_IRE\_Test\_New\_Results.md}.
\bibitem{R4} Data Collection Plan. \emph{IRE-Doc/Data\_Collection\_Plan.md}.
\bibitem{R5} Results Section Draft. \emph{IRE-Doc/Results\_Section\_Draft.md}.
\bibitem{R6} Gaps and Future Works. \emph{IRE-Doc/Gaps \& Future Works.md}.
\bibitem{R7} IEEE Standard 830-1998, \emph{Recommended Practice for Software Requirements Specifications}, IEEE Computer Society.
\bibitem{R8} OWASP Top 10 (2021), owasp.org.
\bibitem{R9} RFC 1918 --- \emph{Address Allocation for Private Internets}, IETF.
\bibitem{R10} \emph{CVE Format Standard}, mitre.org.
\end{thebibliography}

\subsection{Overview}
The remainder of this document is organized as follows:
\begin{itemize}
\item \textbf{Section 2} describes the overall product perspective, user classes, and operating environment;
\item \textbf{Section 3} specifies all system features and functional requirements in detail;
\item \textbf{Section 4} defines external interface requirements (API, hardware, software, communication);
\item \textbf{Section 5} specifies non-functional requirements (performance, security, reliability);
\item \textbf{Section 6} covers data requirements (models, retention, grounding data);
\item \textbf{Section 7} describes verification, validation, and empirical results;
\item \textbf{Section 8} covers deployment procedures;
\item \textbf{Section 9} provides appendices with enumerations, policy tables, and schema definitions.
\end{itemize}
% =================== 2. OVERALL DESCRIPTION ===================
\section{Overall Description}

\subsection{Product Perspective}
The NeuroShell IRE is \textbf{Component 01} of the NeuroShell Framework, an autonomous penetration testing system. The IRE operates as an \textbf{independently deployable microservice} that bridges the gap between human operator intent (natural language) and machine-executable actions (validated JSON contracts).

\begin{lstlisting}
                +------------------------------------+
                |     Operator / Frontend UI         |
                |   (natural-language command)       |
                +----------------+------------------+
                                 | POST /parse
                                 v
                +------------------------------------+
                |     NeuroShell IRE (C01)           |
                |  NLU Gateway + Validation Stack    |
                |  Port 8001 (FastAPI)               |
                +--------+----------------+----------+
                         |                |
                  POST /parse/plan       POST /execute
                         |                |
                +--------v-------+  +----v--------------+
                | Component 02   |  | Component 03      |
                |   Dynamic      |  |  AEERE Executor   |
                |   Planner      |  |  (Port 8003)      |
                | (Port 8002)    |  +-------------------+
                +--------+-------+
                         |
                +--------v-------+
                | Component 04   |
                |    AVAE        |
                |  Vuln Analy.   |
                | (Port 8004)    |
                +----------------+
\end{lstlisting}

The IRE does \textbf{not} depend on downstream components for core operation. Component 02 (Planner), Component 03 (Executor), and Component 04 (Vulnerability Analyzer) are optional integrations controlled by feature flags.

\subsection{Product Functions}
The IRE provides the following core functions:

\begin{table*}[t]
\caption{IRE core product functions.}
\label{tab:functions}
\begin{tabularx}{\textwidth}{l X}
\toprule
\textbf{Function} & \textbf{Description}\\
\midrule
\textbf{Intent Classification} & Classify natural-language commands into one of 9 intent types.\\
\textbf{Target Extraction} & Extract and validate target specification (IP, subnet, domain, URL, hostname).\\
\textbf{Entity Extraction} & Extract ports, modifiers, CVE IDs, tool hints, and scheduling.\\
\textbf{Hallucination Detection} & Detect 6 classes of LLM hallucination via structured validation.\\
\textbf{Adversarial Detection} & Detect prompt injection, jailbreak, and encoded payload attacks.\\
\textbf{RBAC Enforcement} & Enforce role-based access control (4 roles, 9 intents).\\
\textbf{Session Management} & Multi-turn context injection with configurable memory window.\\
\textbf{Semantic Caching} & Exact-hash plus cosine-similarity response cache with grounding invalidation.\\
\textbf{Self-Consistency Voting} & Tier-2 N-sample consensus with semantic entropy uncertainty quantification.\\
\textbf{Contract Generation} & Produce \texttt{PlannerContract} JSON with tool synthesis and command generation.\\
\textbf{Audit Logging} & Privacy-preserving (SHA-256 hashed) append-only audit trail.\\
\textbf{Escalation Queue} & Human-in-the-loop review for high-severity findings with HTTP 202 pattern.\\
\bottomrule
\end{tabularx}
\end{table*}

\subsection{User Classes and Characteristics}
\begin{table*}[t]
\caption{User classes, privileges, and permitted intents.}
\label{tab:roles}
\begin{tabularx}{\textwidth}{l l X l}
\toprule
\textbf{Role} & \textbf{Privilege Level} & \textbf{Permitted Intents} & \textbf{Description}\\
\midrule
\textbf{viewer} & Read-only & None (pre-inference fast-fail) & Observation-only; cannot execute any commands.\\
\textbf{analyst} & Standard & NETWORK\_SCAN, VULNERABILITY\_AUDIT, SERVICE\_ENUMERATION, PASSIVE\_RECON, AMBIGUOUS, REJECTED & Passive and observational operations only.\\
\textbf{operator} & Elevated & All 9 intents & Active operations including exploitation and password attacks.\\
\textbf{admin} & Full & All 9 intents & System administration including model switching and policy management.\\
\bottomrule
\end{tabularx}
\end{table*}

\textbf{Default role:} \texttt{analyst} (when no role is specified in the request).

\textbf{Unknown roles:} Treated as \texttt{viewer} (most restrictive --- implicitly denied everywhere).

\subsection{Operating Environment}

\subsubsection{Hardware}
\begin{table}[h]
\caption{Hardware requirements.}
\label{tab:hardware}
\begin{tabularx}{\columnwidth}{l X X}
\toprule
\textbf{Component} & \textbf{Minimum} & \textbf{Recommended}\\
\midrule
CPU & 4 cores (2.0 GHz+) & 8 cores (3.0 GHz+)\\
RAM & 8 GB & 16 GB\\
GPU & None (CPU-only supported) & NVIDIA CUDA-capable GPU\\
Storage & 10 GB free & 50 GB SSD\\
Network & Loopback (localhost) & LAN (100 Mbps+)\\
\bottomrule
\end{tabularx}
\end{table}

\subsubsection{Software}
\begin{table}[h]
\caption{Software environment.}
\label{tab:software}
\begin{tabularx}{\columnwidth}{l l X}
\toprule
\textbf{Component} & \textbf{Version} & \textbf{Purpose}\\
\midrule
Python & 3.11+ & Runtime\\
Ollama & 0.32.x & LLM inference daemon\\
Docker & 24.x (optional) & Containerized deployment\\
Docker Compose & 2.x (optional) & Multi-service orchestration\\
\bottomrule
\end{tabularx}
\end{table}

\subsubsection{LLM Models}
\begin{table}[h]
\caption{LLM models.}
\label{tab:models}
\begin{tabularx}{\columnwidth}{l l X}
\toprule
\textbf{Model} & \textbf{Size} & \textbf{Purpose}\\
\midrule
\texttt{gemma4:latest} & 9.6 GB & Primary inference (default)\\
\texttt{qwen2.5-coder:7b} & 4.7 GB & Alternative inference backbone\\
\bottomrule
\end{tabularx}
\end{table}

\subsection{Design and Implementation Constraints}
\begin{enumerate}
\item \textbf{Language:} Python 3.11+ (async-first with FastAPI/uvicorn).
\item \textbf{Schema Validation:} Pydantic v2 with \texttt{extra="forbid"} (strict schema enforcement).
\item \textbf{Inference Runtime:} Ollama (local daemon, no cloud API calls).
\item \textbf{Async Concurrency:} Single-worker uvicorn (pipeline is CPU-bound during LLM inference).
\item \textbf{Rate Limiting:} slowapi, 30 requests/minute per IP.
\item \textbf{Input Limit:} 2000 characters per command.
\item \textbf{Session Limit:} 1024 concurrent sessions, 8 turns per session, 30-minute TTL.
\item \textbf{Cache Limits:} 256 entries (inference LRU), 512 entries (semantic cache).
\item \textbf{Privacy:} Raw input never stored in plaintext; SHA-256 hashing only.
\end{enumerate}

\subsection{Assumptions and Dependencies}

\subsubsection{Assumptions}
\begin{enumerate}
\item Ollama daemon is running and accessible at the configured URL.
\item Required LLM models are pre-pulled and available.
\item The operator has valid credentials for the assigned RBAC role.
\item Network connectivity exists between IRE and the Ollama daemon.
\item The engagement scope CIDR is correctly configured for the target environment.
\end{enumerate}

\subsubsection{Dependencies}
\begin{table}[h]
\caption{External dependencies.}
\label{tab:deps}
\begin{tabularx}{\columnwidth}{l l X}
\toprule
\textbf{Dependency} & \textbf{Type} & \textbf{Purpose}\\
\midrule
Ollama & Runtime & LLM inference engine\\
\texttt{data/cve\_grounding.json} & Data & CVE reference dataset (20 entries)\\
\texttt{data/alias\_map.json} & Data & Exploit-name to CVE alias mappings (17 entries)\\
\texttt{all-MiniLM-L6-v2} & ML Model & Sentence embeddings for semantic cache\\
\texttt{en\_core\_web\_sm} & ML Model & spaCy English NLP model\\
\bottomrule
\end{tabularx}
\end{table}
% =================== 3. SYSTEM FEATURES ===================
\section{System Features and Functional Requirements}

\subsection{Pipeline Architecture}

\subsubsection{Overview}
The IRE implements a \textbf{7-stage pipeline with 11 middleware components}, organized as follows:
\begin{lstlisting}
PRE-INFERENCE MIDDLEWARE (M1-M4)
    v
7-STAGE PIPELINE (S1-S7)
    v
POST-INFERENCE MIDDLEWARE (M5-M11)
    v
RESPONSE ASSEMBLY + AUDIT
\end{lstlisting}

\reqid{Requirement FR-PIPE-001:} The pipeline shall process each request through the full stage sequence (S1$\rightarrow$S7) unless an earlier stage raises a blocking exception.

\reqid{Requirement FR-PIPE-002:} Each middleware and stage shall be wrapped in individual try/except blocks so that a failure in non-critical middleware degrades gracefully rather than aborting the request.

\reqid{Requirement FR-PIPE-003:} The pipeline shall accumulate \texttt{ValidationFindings} across all stages (S5/S6/S6B/S7) and include them in the response payload regardless of pass/fail status.

\reqid{Requirement FR-PIPE-004:} Every non-passed finding shall be written to the SQLite audit log via \texttt{\_log\_audit\_findings}.

\subsubsection{Stage Sequence}
\begin{table*}[t]
\caption{Pipeline stage sequence (S1--S7).}
\label{tab:stages}
\begin{tabularx}{\textwidth}{l l X X X}
\toprule
\textbf{Stage} & \textbf{Module} & \textbf{Function} & \textbf{Input} & \textbf{Output}\\
\midrule
S1 & \texttt{InputNormalizer} & NFKC Unicode normalization, control-char strip, jargon standardization & Raw text & Normalized text\\
S2 & \texttt{AliasResolver} & Alias to CVE canonical mapping with provenance annotation & Normalized text & Enriched text\\
S3 & \texttt{OllamaInferenceEngine} & Tier-1/Tier-2 LLM inference with self-consistency & Enriched text & \texttt{IntentSchema}\\
S4 & \texttt{JSONParser} & Strip thought tags, fences; extract JSON object & LLM output text & Parsed dict\\
S5 & \texttt{SchemaValidator} & Pydantic v2 \texttt{IntentSchema} validation & Dict & \texttt{IntentSchema}\\
S6 & \texttt{RegexValidator} & IP/CIDR/CVE/port/metachar regex checks & \texttt{IntentSchema} & Validated \texttt{IntentSchema}\\
S6B & \texttt{NetworkArchitectureValidator} & Type consistency, \texttt{ipaddress} parse, scope containment, CIDR sanity & \texttt{IntentSchema} & Validated \texttt{IntentSchema}\\
S7 & \texttt{ScopeGuard} & 7 injection regexes on raw input plus RFC-1918 private-range check & Raw input + \texttt{IntentSchema} & (\texttt{IntentSchema}, warnings)\\
\bottomrule
\end{tabularx}
\end{table*}

\subsubsection{Middleware Sequence}
\begin{table*}[t]
\caption{Middleware sequence (M1--M11).}
\label{tab:middleware}
\begin{tabularx}{\textwidth}{l l l X}
\toprule
\textbf{Middleware} & \textbf{Module} & \textbf{Phase} & \textbf{Function}\\
\midrule
M1 & \texttt{AdversarialDetector} & Pre-inference & Weighted threat scoring $\rightarrow$ block/allow\\
M2 & \texttt{RBACGuard.pre\_inference\_check} & Pre-inference & Restricted-role fast-fail\\
M3 & \texttt{SessionContextStore.inject\_context} & Pre-inference & Prepend session turn history\\
M4 & \texttt{SemanticCache.lookup} & Pre-inference & Exact-hash then cosine cache hit\\
M5 & \texttt{RBACGuard.post\_inference\_check} & Post-inference & Role $\times$ intent authorization\\
M7 & \texttt{SemanticCache.store} & Post-inference & Store successful response\\
M8 & \texttt{IREResponseV2.from\_v1} & Post-inference & Assemble V2 response with metadata\\
M9 & \texttt{SubIntentClassifier.enrich} & Post-inference & Hierarchical sub-intent classification\\
M10 & \texttt{SessionContextStore.update} & Post-inference & Record completed turn\\
M11 & Session info enrich & Post-inference & Attach turn count to response\\
\bottomrule
\end{tabularx}
\end{table*}

\subsection{Pre-Inference Middleware (M1--M4)}

\subsubsection{M1 --- AdversarialDetector}
\reqid{Requirement FR-ADV-001:} The adversarial detector shall scan raw input against 15 lexical injection patterns (weights 0.8--1.0) and 7 structural anomaly patterns (weights 0.5--0.9) plus 3 heuristic checks.

\reqid{Requirement FR-ADV-002:} Any single signal with weight $\geq$ 0.95 shall trigger immediate blocking (auto-block), regardless of cumulative score.

\reqid{Requirement FR-ADV-003:} The cumulative threat score shall be capped at 1.0, and blocking shall occur when the score meets or exceeds the configured threshold (default 0.70).

\reqid{Requirement FR-ADV-004:} Scanning shall short-circuit as soon as the threshold is reached (efficiency optimization).

\reqid{Requirement FR-ADV-005:} The adversarial detector may be disabled via the \texttt{adversarial\_detection} feature flag.

Detection categories:
\begin{table}[h]
\caption{Adversarial detection layers.}
\label{tab:advlayers}
\begin{tabularx}{\columnwidth}{l >{\raggedright\arraybackslash}p{0.18\columnwidth} X >{\raggedright\arraybackslash}p{0.3\columnwidth}}
\toprule
\textbf{Layer} & \textbf{Patterns} & \textbf{Weight Range} & \textbf{Examples}\\
\midrule
A: Lexical & 15 regex patterns & 0.8--1.0 & \texttt{ignore previous instructions}, \texttt{jailbreak}, \texttt{$<$|im\_start|$>$}\\
B: Structural & 7 regex patterns & 0.5--0.9 & Hex encoding, base64 payloads, \texttt{$<$script$>$}, \texttt{eval()}\\
B: Heuristic & 3 checks & 0.3--0.4 & Special-char density $>$30\%, token $>$100 chars, repeated n-grams\\
\bottomrule
\end{tabularx}
\end{table}

Decision rule:
\begin{equation}
\mathrm{BLOCK} \iff \bigl(\exists\, i \in M(x):\ w_{i} \geq 0.95\bigr) \lor \bigl(S(x) \geq \tau\bigr),
\end{equation}
\begin{equation}
S(x) = \min\Bigl(1,\ \sum_{i \in M(x)} w_{i}\Bigr), \qquad \tau = 0.70.
\end{equation}

\subsubsection{M2 --- RBACGuard (Pre-Inference)}
\reqid{Requirement FR-RBAC-001:} The pre-inference check shall fast-fail for any role in the \texttt{RESTRICTED\_ROLES} set (\texttt{viewer}).

\reqid{Requirement FR-RBAC-002:} An \texttt{RBACError} shall be raised with \texttt{intent="ANY"} for restricted roles.

\reqid{Requirement FR-RBAC-003:} The pre-check may be disabled via the \texttt{rbac} feature flag.

\subsubsection{M3 --- SessionContextStore}
\reqid{Requirement FR-SES-001:} The session context store shall maintain per-session turn logs with configurable limits: max 8 turns, max 1024 sessions, 30-minute TTL, LRU eviction.

\reqid{Requirement FR-SES-002:} Context shall be injected in the following format:
\begin{lstlisting}
Session context:
[Turn 1] NETWORK_SCAN on SUBNET "192.168.1.0/24" ports=[22,80]
[Turn 2] VULNERABILITY_AUDIT on IP "10.0.0.5" cves=[CVE-2021-44228]
Current command: <user command>
\end{lstlisting}

\reqid{Requirement FR-SES-003:} Context injection shall occur BEFORE input normalization and BEFORE cache lookup, so cache keys include session state.

\reqid{Requirement FR-SES-004:} Only successful responses shall be recorded as turns (error responses are not stored).

\subsubsection{M4 --- SemanticCache}
\reqid{Requirement FR-CACHE-001:} The semantic cache shall support two-tier lookup: exact SHA-256 hash match first, then cosine similarity over all entries.

\reqid{Requirement FR-CACHE-002:} The cosine similarity threshold shall be configurable (default 0.82).

\reqid{Requirement FR-CACHE-003:} Each cache entry shall carry a \texttt{grounding\_hash} (SHA-256 of \texttt{cve\_grounding.json} + \texttt{alias\_map.json}). Entries with stale grounding hashes shall be purged on contact.

\reqid{Requirement FR-CACHE-004:} The cache shall use LRU eviction at capacity (default 512 entries).

\reqid{Requirement FR-CACHE-005:} Only successful responses shall be cached.

\reqid{Requirement FR-CACHE-006:} The embedding model (\texttt{all-MiniLM-L6-v2}) shall be lazy-loaded on first use.

\subsection{Core Pipeline Stages (S1--S7)}

\subsubsection{S1 --- InputNormalizer}
\reqid{Requirement FR-NORM-001:} Input normalization shall apply: NFKC Unicode normalization, control-char stripping (ord $<$ 32), think-token removal (\texttt{ thinking}, \texttt{$<$|think|$>$}), whitespace collapse.

\reqid{Requirement FR-NORM-002:} A 20-term jargon standardization table shall map offensive-security slang to canonical terms (e.g., \texttt{syn scan} $\rightarrow$ \texttt{SYN stealth scan}, \texttt{pwn} $\rightarrow$ \texttt{exploit}, \texttt{privesc} $\rightarrow$ \texttt{privilege escalation}).

\reqid{Requirement FR-NORM-003:} Empty or whitespace-only input shall raise an \texttt{InputValidationError} with an HTTP 200 error envelope.

\subsubsection{S2 --- AliasResolver}
\reqid{Requirement FR-ALIAS-001:} The alias resolver shall load \texttt{data/alias\_map.json} (17 active entries) and compile word-boundary patterns sorted longest-first.

\reqid{Requirement FR-ALIAS-002:} Aliases shall be rewritten to \texttt{CANONICAL (alias: ORIGINAL)} format so the LLM sees provenance.

\reqid{Requirement FR-ALIAS-003:} Each alias shall be rewritten only at its first occurrence per input.

\subsubsection{S3 --- OllamaInferenceEngine (Tier-1/Tier-2)}
\reqid{Requirement FR-INF-001:} The inference engine shall support two tiers of inference with conditional gating.

\reqid{Requirement FR-INF-002:} Tier-1 shall execute a single Ollama API call at the configured temperature (default 0.1).

\reqid{Requirement FR-INF-003:} The Tier-1 confidence proxy shall be computed as
\begin{equation}
c_{\mathrm{base}} = \begin{cases} o.\mathrm{confidence}, & \text{if parsed to } [0,1],\\ 0.5, & \text{otherwise.} \end{cases}
\end{equation}
\begin{equation}
c_{1} = \begin{cases} \min(c_{\mathrm{base}},\,0.4), & i \in \mathcal{H},\\ \min(c_{\mathrm{base}},\,0.5), & i \notin \mathcal{H},\ t_{v} \text{ missing},\\ c_{\mathrm{base}}, & \text{otherwise,} \end{cases}
\end{equation}
where $i = \mathrm{intent}$, $t_{v} = \mathrm{target.value}$, and $\mathcal{H} = \{\mathrm{AMBIGUOUS},\ \mathrm{REJECTED},\ \emptyset\}$.

\reqid{Requirement FR-INF-004:} If $c_{1} \geq \gamma$ ($\gamma = 0.7$), the Tier-1 output shall be returned immediately (entropy logged as 0, band ``low'', tier ``1'').

\reqid{Requirement FR-INF-005:} Tier-2 shall execute when $c_{1} < \gamma$, generating $N=5$ total samples (Tier-1 reused $+$ $N-1$ resamples) at a higher temperature (default 0.7).

\reqid{Requirement FR-INF-006:} Tier-2 outputs shall be clustered by semantic-equivalence signature
\begin{equation}
\begin{aligned}
\sigma(o) ={}& \bigl(\mathrm{intent},\ \mathrm{target\_type},\ \mathrm{target\_value.lower()},\\
&\quad \mathrm{tuple}(\mathrm{sorted}(cve\_ids))\bigr),
\end{aligned}
\end{equation}
which deliberately ignores ports and modifiers.

\reqid{Requirement FR-INF-007:} Semantic entropy shall be computed as
\begin{equation}
H(x) = -\sum_{j} p_{j}\, \ln p_{j}, \qquad p_{j} = |C_{j}|/N.
\end{equation}

\reqid{Requirement FR-INF-008:} Entropy bands shall be mapped via configurable thresholds:
\begin{itemize}
\item $H < 0.3$ $\rightarrow$ ``low'';
\item $0.3 \leq H \leq 0.8$ $\rightarrow$ ``medium'';
\item $H > 0.8$ $\rightarrow$ ``high''.
\end{itemize}

\reqid{Requirement FR-INF-009:} The resolution method for Tier-2 shall be plurality vote (largest cluster wins; ties resolve by insertion order).

\reqid{Requirement FR-INF-010:} The inference engine shall maintain a two-layer cache: in-memory LRU (256 entries) plus JSON disk persistence, keyed by \texttt{SHA-256(model\_name|enriched\_input)}.

\subsubsection{S4 --- JSONParser}
\reqid{Requirement FR-JSON-001:} The JSON parser shall strip \texttt{ thinking...response}, \texttt{$<$|think|$>$}, and markdown code fences.

\reqid{Requirement FR-JSON-002:} It shall extract the substring from the first \texttt{\{} to the last \texttt{\}} (inclusive).

\reqid{Requirement FR-JSON-003:} It shall validate that the result is a Python dict; failure raises \texttt{JSONParseError}.

\subsubsection{S5 --- SchemaValidator}
\reqid{Requirement FR-SCH-001:} The schema validator shall validate the parsed dict against the \texttt{IntentSchema} Pydantic v2 model via \texttt{model\_validate()}.

\reqid{Requirement FR-SCH-002:} On validation failure, a \texttt{ValidationFinding} shall be created with \texttt{severity="block"} and the appropriate \texttt{HallucinationClass} mapped from the failing field:
\begin{table*}[t]
\caption{Schema error to hallucination-class mapping.}
\label{tab:schmap}
\begin{tabularx}{\textwidth}{>{\ttfamily}l >{\ttfamily}X}
\toprule
\textbf{Error Field/Condition} & \textbf{Hallucination Class}\\
\midrule
\texttt{ports} in error path & FABRICATED\_PARAMETER\\
\texttt{cve\_ids} in error path & FABRICATED\_CVE\\
\texttt{target} or \texttt{intent} in error path & TARGET\_TYPE\_MISMATCH\\
\texttt{rejection\_reason}, \texttt{confidence}, AMBIGUOUS, REJECTED in message & UNGROUNDED\_CONFIDENCE\\
Default fallback & UNGROUNDED\_CONFIDENCE\\
\bottomrule
\end{tabularx}
\end{table*}

\reqid{Requirement FR-SCH-003:} Schema constraints include:
\begin{itemize}
\item \texttt{intent} $\in$ 9-class \texttt{IntentType} enum;
\item \texttt{target.type} $\in$ 6-class \texttt{TargetType} enum;
\item \texttt{confidence} $\in$ [0.0,~1.0];
\item \texttt{ports[$i$]} $\in$ [1,~65535];
\item \texttt{cve\_ids[$i$]} matches \texttt{\^{}CVE-\textbackslash d\{4\}-\textbackslash d\{4,7\}\$};
\item \texttt{intent=REJECTED} $\Rightarrow$ \texttt{rejection\_reason} required;
\item \texttt{intent=AMBIGUOUS} $\Rightarrow$ \texttt{confidence} $<$ 0.5;
\item No extra fields allowed (\texttt{extra="forbid"}).
\end{itemize}

\subsubsection{S6 --- RegexValidator}
\reqid{Requirement FR-REG-001:} Four sequential checks shall be performed:
\begin{table*}[t]
\caption{RegexValidator checks.}
\label{tab:regex}
\begin{tabularx}{\textwidth}{l X l l}
\toprule
\textbf{Check} & \textbf{Description} & \textbf{Failure Class} & \textbf{Severity}\\
\midrule
1. Target Format & Type-specific regex: IPv4, IPv6, CIDR, Domain, Hostname, URL & TARGET\_TYPE\_MISMATCH & block\\
2. CVE Syntax + Grounding & \texttt{\^{}CVE-\textbackslash d\{4\}-\textbackslash d\{4,7\}\$} syntax; then lookup in \texttt{cve\_grounding.json} & FABRICATED\_CVE / UNGROUNDED\_CONFIDENCE & block / from policy\\
3. Port Bounds & $\forall p$: $1 \leq p \leq 65535$ & FABRICATED\_PARAMETER & block\\
4. Shell Metacharacters & Scan for \texttt{[;\&|`\$!$>$]} in \texttt{target.value} and modifiers & FABRICATED\_PARAMETER & block\\
\bottomrule
\end{tabularx}
\end{table*}

\reqid{Requirement FR-REG-002:} CVE syntax failures shall produce severity from \texttt{EnforcementPolicy} (default: block).

\reqid{Requirement FR-REG-003:} Syntactically valid but ungrounded CVEs (not in \texttt{cve\_grounding.json}) shall produce \texttt{UNGROUNDED\_CONFIDENCE} with severity from policy (default: warn --- advisory only).

\subsubsection{S6B --- NetworkArchitectureValidator}
\reqid{Requirement FR-NET-001:} Four-step validation for IP/SUBNET targets:
\begin{table*}[t]
\caption{Network architecture validation steps.}
\label{tab:net}
\begin{tabularx}{\textwidth}{l X >{\raggedright\arraybackslash}p{0.2\textwidth}}
\toprule
\textbf{Step} & \textbf{Check} & \textbf{Failure Class}\\
\midrule
1. Type Consistency & IP must not contain ``/''; SUBNET must contain ``/''; UNKNOWN must be empty & TARGET\_TYPE\_MISMATCH\\
2. Parseability & \texttt{ipaddress.ip\_network()} or \texttt{ip\_address()} parse & format error\\
3. Scope Containment & IP: $\mathrm{addr} \in \Omega$; SUBNET: $\mathrm{net} \subseteq \Omega$ (\texttt{subnet\_of}) & OUT\_OF\_SCOPE\_TARGET\\
4. CIDR Sanity & Prefix $\geq$ \texttt{max\_cidr\_prefix} (30); reject /32 subnet; IPv6 $<$ /64 rejected & OUT\_OF\_SCOPE\_TARGET\\
\bottomrule
\end{tabularx}
\end{table*}

\reqid{Requirement FR-NET-002:} Domain, URL, Hostname, and UNKNOWN types shall auto-pass the scope check (scope validation not applicable).

\reqid{Requirement FR-NET-003:} Engagement mode \texttt{"disabled"} shall short-circuit all validation; \texttt{"warn"} shall downgrade all severities to advisory warnings.

\subsubsection{S7 --- ScopeGuard}
\reqid{Requirement FR-SCOPE-001:} The scope guard shall scan the \textbf{raw input} (not the enriched/normalized form) against 7 injection regex patterns. Any match shall produce an unconditional hard block.

\reqid{Requirement FR-SCOPE-002:} Private-range membership shall be checked against 7 network ranges:
\begin{itemize}
\item \texttt{10.0.0.0/8}, \texttt{172.16.0.0/12}, \texttt{192.168.0.0/16};
\item \texttt{127.0.0.0/8} (loopback), \texttt{169.254.0.0/16} (link-local);
\item \texttt{::1/128} (IPv6 loopback), \texttt{fc00::/7} (IPv6 ULA).
\end{itemize}

\reqid{Requirement FR-SCOPE-003:} Public (non-private) targets shall trigger \texttt{OUT\_OF\_SCOPE\_TARGET} with severity determined by \texttt{EnforcementPolicy}:
\begin{itemize}
\item \texttt{"strict"} mode $\rightarrow$ block;
\item \texttt{"research"} mode $\rightarrow$ warn (advisory).
\end{itemize}

\subsection{Post-Inference Middleware (M5--M11)}

\subsubsection{M5 --- RBACGuard (Post-Inference)}
\reqid{Requirement FR-RBAC-010:} The post-inference check shall verify that the detected intent is permitted for the caller's role using \texttt{INTENT\_ROLE\_POLICY}.

\reqid{Requirement FR-RBAC-011:} Denied intents shall raise \texttt{RBACError} listing the detected intent and the role's permitted intents.

\subsubsection{M7 --- SemanticCache (Store)}
\reqid{Requirement FR-CACHE-010:} Successful responses shall be stored with their grounding hash for future invalidation.

\subsubsection{M8 --- IREResponseV2 Assembly}
\reqid{Requirement FR-RESP-001:} The V2 response shall include all V1 fields plus: \texttt{sub\_intent}, \texttt{secondary\_intents}, \texttt{uncertainty\_band}, \texttt{raw\_entropy}, \texttt{tier\_triggered}, \texttt{resolution\_method}, \texttt{cache\_hit}, \texttt{rbac\_role}, \texttt{session\_turns}, \texttt{validation\_findings}, \texttt{schema\_version}.

\reqid{Requirement FR-RESP-002:} Error responses shall follow the 200-envelope pattern: HTTP 200 with \texttt{status:"error"}, carrying \texttt{error}, \texttt{stage}, \texttt{detail}, \texttt{field}, \texttt{validation\_findings}, and \texttt{latency\_ms}.

\subsubsection{M9 --- SubIntentClassifier}
\reqid{Requirement FR-SUB-001:} The sub-intent classifier shall be a first-match-wins rule table over 27 sub-intent types across 7 parent intents.

\reqid{Requirement FR-SUB-002:} Classification shall be gated: disabled for AMBIGUOUS/REJECTED intents and for confidence $<$ 0.6.

\subsubsection{M10/M11 --- Session Update}
\reqid{Requirement FR-SES-010:} After successful processing, the session store shall record the turn with intent, target, ports, CVE IDs, and confidence.
\subsection{Schema and Contract System}

\subsubsection{IntentSchema (IRE Internal)}
The \texttt{IntentSchema} is the IRE's internal validated representation:
\begin{table*}[t]
\caption{IntentSchema fields.}
\label{tab:intentschema}
\begin{tabularx}{\textwidth}{>{\ttfamily}l l l X}
\toprule
\textbf{Field} & \textbf{Type} & \textbf{Required} & \textbf{Constraints}\\
\midrule
\texttt{intent} & IntentType & Yes & 9-class enum\\
\texttt{target} & Target & Yes & \texttt{\{type: TargetType, value: str\}}\\
\texttt{ports} & List[int] & No & Each $\in$ [1, 65535]\\
\texttt{modifiers} & List[str] & No & Free text\\
\texttt{cve\_ids} & List[str] & No & Each matches \texttt{\^{}CVE-\textbackslash d\{4\}-\textbackslash d\{4,7\}\$}\\
\texttt{tool\_hint} & Optional[str] & No & Suggested tool\\
\texttt{schedule} & Optional[str] & No & Cron expression\\
\texttt{confidence} & float & Yes & $\in$ [0.0, 1.0]\\
\texttt{rejection\_reason} & Optional[str] & Yes if REJECTED & Required when \texttt{intent=REJECTED}\\
\bottomrule
\end{tabularx}
\end{table*}

\subsubsection{PlannerContract (Component 02 Hand-off)}
The \texttt{PlannerContract} extends the intent schema with downstream-consumption fields:
\begin{table}[h]
\caption{PlannerContract fields.}
\label{tab:plannercontract}
\begin{tabularx}{\columnwidth}{>{\ttfamily}l l X}
\toprule
\textbf{Field} & \textbf{Type} & \textbf{Purpose}\\
\midrule
\texttt{sub\_intent} & Optional[SubIntentType] & Hierarchical refinement for C02 command construction\\
\texttt{tool\_hint} & Optional[str] & Tool suggestion (nmap, gobuster, nikto, manual, none)\\
\texttt{schedule} & Optional[str] & Cron expression; None = execute immediately\\
\texttt{scope\_warnings} & List[str] & SCOPE\_WARNING and ARCH\_WARNING entries\\
\bottomrule
\end{tabularx}
\end{table}

Executability predicates:
\begin{itemize}
\item \texttt{is\_executable} $\rightarrow$ False if REJECTED/AMBIGUOUS, or if any \texttt{scope\_warning} contains ``SCOPE\_VIOLATION'' or ``ARCH\_WARNING'';
\item \texttt{is\_immediate} $\rightarrow$ True if \texttt{schedule} is None;
\item \texttt{has\_scope\_issues} $\rightarrow$ True if any \texttt{scope\_warnings} present.
\end{itemize}

\subsection{Hallucination Taxonomy and Enforcement}

\subsubsection{HallucinationClass Enum (6 Classes)}
\begin{table*}[t]
\caption{Hallucination classes.}
\label{tab:hallucination}
\begin{tabularx}{\textwidth}{>{\ttfamily}l X X}
\toprule
\textbf{Class} & \textbf{Description} & \textbf{Detection Stage}\\
\midrule
FABRICATED\_CVE & Malformed or non-existent CVE ID & S5 (schema), S6 (regex+grounding)\\
TARGET\_TYPE\_MISMATCH & Structural mismatch between target type and value & S5 (schema), S6 (format), S6B (consistency)\\
CONTRADICTORY\_ACTION\_TARGET & Incompatible action-target pairing & S5 (schema)\\
FABRICATED\_PARAMETER & Out-of-bounds or injected parameters & S5 (schema), S6 (port/metachar)\\
UNGROUNDED\_CONFIDENCE & Valid syntax but unindexed CVE; low-confidence output & S6 (grounding), S5 (confidence)\\
OUT\_OF\_SCOPE\_TARGET & Target outside authorized engagement scope & S6B (scope), S7 (RFC-1918)\\
\bottomrule
\end{tabularx}
\end{table*}

\subsubsection{Enforcement Policy}
\begin{table}[h]
\caption{Default enforcement by class.}
\label{tab:enforcement}
\begin{tabularx}{\columnwidth}{>{\ttfamily}l l X}
\toprule
\textbf{HallucinationClass} & \textbf{Default Action} & \textbf{Reasoning}\\
\midrule
FABRICATED\_PARAMETER & BLOCK & Execution risk from out-of-bounds parameters or injection\\
CONTRADICTORY\_ACTION\_TARGET & BLOCK & Safety invariant violation\\
TARGET\_TYPE\_MISMATCH & BLOCK & Structural mismatch prevents downstream execution\\
FABRICATED\_CVE & BLOCK & Complete LLM hallucination\\
UNGROUNDED\_CONFIDENCE & WARN & Novel but plausible CVEs; advisory only\\
OUT\_OF\_SCOPE\_TARGET & Dynamic & BLOCK in strict mode; WARN in research mode\\
\bottomrule
\end{tabularx}
\end{table}

\subsubsection{ValidationFinding Schema}
Each validation check produces a \texttt{ValidationFinding}:
\begin{table*}[t]
\caption{ValidationFinding fields.}
\label{tab:vf}
\begin{tabularx}{\textwidth}{>{\ttfamily}l l X}
\toprule
\textbf{Field} & \textbf{Type} & \textbf{Values}\\
\midrule
\texttt{validator} & str & Which component performed the check\\
\texttt{passed} & bool & True if check passed\\
\texttt{hallucination\_class} & Optional[HallucinationClass] & Assigned class on failure\\
\texttt{detail} & str & Human-readable description\\
\texttt{severity} & Literal["block","warn","escalate"] & Enforcement action\\
\bottomrule
\end{tabularx}
\end{table*}

\subsection{Contract Builder and Tool Synthesis}

\subsubsection{Tool Synthesis Table}
\begin{table*}[t]
\caption{Tool synthesis per intent.}
\label{tab:toolsynth}
\begin{tabularx}{\textwidth}{>{\ttfamily}l l X}
\toprule
\textbf{Intent} & \textbf{Tool} & \textbf{Command Template}\\
\midrule
NETWORK\_SCAN & nmap & \texttt{nmap [flags] $<$target$>$}\\
VULNERABILITY\_AUDIT & nikto (URL) / nmap & \texttt{nikto -h $<$target$>$ [-p port] [-ssl]}\\
SERVICE\_ENUMERATION & nmap & \texttt{nmap -sV [flags] $<$target$>$}\\
EXPLOITATION & manual & Requires human authorization\\
PASSWORD\_ATTACK & manual & Lockout risk warning\\
DIRECTORY\_BRUTEFORCE & gobuster & \texttt{gobuster dir -u $<$url$>$ -w $<$wordlist$>$ -q}\\
PASSIVE\_RECON & none & Information gathering only\\
AMBIGUOUS & none & Requires clarification\\
REJECTED & none & Command rejected\\
\bottomrule
\end{tabularx}
\end{table*}

\subsubsection{Priority Ladder}
\begin{table}[h]
\caption{Contract priority ladder.}
\label{tab:priority}
\begin{tabularx}{\columnwidth}{l X}
\toprule
\textbf{Priority} & \textbf{Condition}\\
\midrule
SCHEDULED & Cron expression is set\\
CRITICAL & Intent is EXPLOITATION\\
HIGH & Intent is VULNERABILITY\_AUDIT with CVE IDs\\
MEDIUM & NETWORK\_SCAN, SERVICE\_ENUMERATION, or VULNERABILITY\_AUDIT (no CVEs)\\
LOW & All other intents\\
\bottomrule
\end{tabularx}
\end{table}

\subsubsection{Nmap Flag Mapping}
\begin{table}[h]
\caption{Modifier to nmap-flag mapping.}
\label{tab:nmapflags}
\begin{tabularx}{\columnwidth}{>{\ttfamily}l >{\ttfamily}l}
\toprule
\textbf{Modifier} & \textbf{Flag}\\
\midrule
\texttt{stealth} & \texttt{-sS} (SYN scan)\\
\texttt{aggressive} & \texttt{-A}\\
\texttt{version\_detect} & \texttt{-sV}\\
\texttt{os\_detect} & \texttt{-O}\\
\texttt{udp} & \texttt{-sU}\\
\texttt{full\_port} & \texttt{-p-}\\
\texttt{verbose} & \texttt{-v}\\
\texttt{stealth} timing & \texttt{-T1}\\
Default timing & \texttt{-T3}\\
\bottomrule
\end{tabularx}
\end{table}

\subsection{Audit and Escalation System}

\subsubsection{Audit Logger}
\reqid{Requirement FR-AUDIT-001:} The audit logger shall maintain an append-only SQLite \texttt{audit\_log} table.

\reqid{Requirement FR-AUDIT-002:} Raw user input shall NEVER be stored in plaintext; only SHA-256 hashes shall be recorded.

\reqid{Requirement FR-AUDIT-003:} Each finding shall be recorded with: timestamp, session\_id, raw\_input\_hash, intent\_summary, hallucination\_class, severity, enforcement\_action, and status.

\subsubsection{Escalation Queue}
\reqid{Requirement FR-ESC-001:} Findings with \texttt{severity="escalate"} shall be recorded with \texttt{status="pending"} in the audit log.

\reqid{Requirement FR-ESC-002:} The \texttt{POST /parse/plan} endpoint shall return HTTP 202 with \texttt{escalation\_id} when an ESCALATE-severity finding is encountered, withholding the contract.

\reqid{Requirement FR-ESC-003:} Admins may approve or reject escalations via \texttt{POST /admin/escalations/\{id\}/resolve}, which releases or discards the cached contract.
% =================== 4. EXTERNAL INTERFACES ===================
\section{External Interface Requirements}

\subsection{User Interfaces}
The IRE does not have a direct user interface. It provides a REST API consumed by:
\begin{itemize}
\item Frontend UI (separate project, \texttt{apps/web/});
\item Downstream components (C02, C03, C04);
\item CLI tools and scripts.
\end{itemize}

\subsection{Hardware Interfaces}
\begin{table}[h]
\caption{Hardware interfaces.}
\label{tab:hwinterface}
\begin{tabularx}{\columnwidth}{l X}
\toprule
\textbf{Interface} & \textbf{Specification}\\
\midrule
Network & TCP/IP over Ethernet or loopback\\
Ollama Communication & HTTP/1.1 to Ollama daemon (default port 11434)\\
Storage & Filesystem for SQLite audit log, JSON cache files, grounding data\\
\bottomrule
\end{tabularx}
\end{table}

\subsection{Software Interfaces}
\begin{table*}[t]
\caption{Software interface map.}
\label{tab:swinterfaces}
\begin{tabularx}{\textwidth}{l l X}
\toprule
\textbf{Interface} & \textbf{Protocol} & \textbf{Target}\\
\midrule
Ollama API & HTTP POST \texttt{/api/generate} & Local Ollama daemon\\
Component 02 & HTTP POST \texttt{/plan} & Dynamic Planner (port 8002, optional)\\
Component 03 & HTTP POST & AEERE Executor (port 8003, optional)\\
Component 04 & HTTP POST & AVAE Vulnerability Analyzer (port 8004, optional)\\
SQLite & File I/O & \texttt{data/audit\_log.db}\\
Filesystem & Read & \texttt{data/cve\_grounding.json}, \texttt{data/alias\_map.json}\\
\bottomrule
\end{tabularx}
\end{table*}

\subsection{Communication Interfaces}

\subsubsection{REST API Base}
\begin{itemize}
\item \textbf{Protocol:} HTTP/1.1 (HTTPS recommended for production);
\item \textbf{Content-Type:} \texttt{application/json};
\item \textbf{Authentication:} \texttt{X-API-Key} header (Bearer-style);
\item \textbf{Rate Limiting:} 30 requests/minute per IP (slowapi);
\item \textbf{CORS:} Allow all origins (development mode).
\end{itemize}

\subsubsection{Response Headers}
\begin{table}[h]
\caption{IRE response headers.}
\label{tab:headers}
\begin{tabularx}{\columnwidth}{>{\ttfamily}l X}
\toprule
\textbf{Header} & \textbf{Description}\\
\midrule
\texttt{X-IRE-Schema-Version} & Schema version(s) supported: \texttt{"1,2"}\\
\texttt{X-IRE-Latency-Ms} & Pipeline processing latency in milliseconds\\
\texttt{X-IRE-Intent} & Detected intent type\\
\texttt{X-IRE-Cache-Hit} & Cache hit type: \texttt{"exact"}, \texttt{"semantic"}, or \texttt{"none"}\\
\texttt{X-Chain-Cache-Hit} & Chain cache hit (on \texttt{/execute} only)\\
\bottomrule
\end{tabularx}
\end{table}

\subsection{API Endpoint Specifications}

\subsubsection{GET /}
\textbf{Purpose:} Service information (no auth required).

\textbf{Response 200:}
\begin{lstlisting}
{
  "service": "NeuroShell IRE",
  "version": "2.0.0",
  "status": "running",
  "docs": "/docs",
  "endpoints": { ... }
}
\end{lstlisting}

\subsubsection{GET /health}
\textbf{Purpose:} System health check (no auth required).

\textbf{Response 200:}
\begin{lstlisting}
{
  "status": "healthy",
  "ollama": { "connected": true, "model": "gemma4:latest" },
  "pipeline": { "stages": 7, "features": { ... } },
  "sessions": { "active": 0, "max": 1024 },
  "semantic_cache": { "size": 0, "max_size": 512 },
  "inference_cache": { "size": 0 }
}
\end{lstlisting}

\textbf{Response 503:} Service unavailable.

\subsubsection{POST /parse}
\textbf{Purpose:} Core intent recognition endpoint.

\textbf{Request:}
\begin{lstlisting}
{
  "command": "scan 192.168.1.0/24 for open ports",
  "session_id": "optional-session-id",
  "role": "analyst"
}
\end{lstlisting}

\textbf{Constraints:}
\begin{itemize}
\item \texttt{command}: required, non-empty, $\leq$ 2000 characters;
\item \texttt{role}: optional, default \texttt{"analyst"}.
\end{itemize}

\textbf{Response 200 (success):}
\begin{lstlisting}
{
  "status": "success",
  "intent": "NETWORK_SCAN",
  "target": { "type": "SUBNET", "value": "192.168.1.0/24" },
  "ports": [],
  "modifiers": [],
  "cve_ids": [],
  "tool_hint": null,
  "schedule": null,
  "confidence": 0.85,
  "rejection_reason": null,
  "scope_warnings": [],
  "latency_ms": 27650,
  "sub_intent": "SYN_STEALTH",
  "uncertainty_band": "low",
  "raw_entropy": 0.0,
  "tier_triggered": 1,
  "resolution_method": "single_sample",
  "cache_hit": "none",
  "rbac_role": "analyst",
  "validation_findings": [],
  "schema_version": 2
}
\end{lstlisting}

\textbf{Response 200 (error envelope):}
\begin{lstlisting}
{
  "status": "error",
  "error": "SCOPE_VIOLATION",
  "stage": "scope_guard",
  "detail": "Target outside authorized scope",
  "latency_ms": 150,
  "validation_findings": [...]
}
\end{lstlisting}

\textbf{Other Responses:} 400 (bad request), 401 (unauthorized), 403 (RBAC denied), 413 (payload too large), 429 (rate limited), 503 (service unavailable).

\subsubsection{POST /parse/plan}
\textbf{Purpose:} Intent recognition plus \texttt{PlannerContract} generation for Component 02.

\textbf{Request:} Same as \texttt{/parse}.

\textbf{Response 200:} \texttt{PlannerContract} JSON (see Section 3.5.2).

\textbf{Response 202:} Escalation pending --- returns \texttt{escalation\_id} and \texttt{status: "pending\_review"}, withholds contract.

\textbf{Response 422:} Contract validation failure.

\subsubsection{POST /execute}
\textbf{Purpose:} Full synchronous C1$\rightarrow$C2$\rightarrow$C3$\rightarrow$C4 chain execution.

\textbf{Request:} Same as \texttt{/parse}.

\textbf{Response 200:} Combined multi-component result (IRE parse + planner output + executor output + vulnerability analysis).

\subsubsection{GET /metrics}
\textbf{Purpose:} Pipeline performance metrics (requires auth).

\textbf{Response 200:}
\begin{lstlisting}
{
  "intent_distribution": { ... },
  "error_counts": { ... },
  "latency_percentiles": { "p50": 16498, "p95": 70893, "p99": ... },
  "entropy_records": { ... },
  "cache_hits": { "exact": 0, "semantic": 0 }
}
\end{lstlisting}

\subsubsection{GET/DELETE /session/\{session\_id\}}
\textbf{Purpose:} Session management (requires auth).

\textbf{GET Response 200:} Session info with turns list.

\textbf{DELETE Response 200:}
\begin{lstlisting}
{ "cleared": true }
\end{lstlisting}

\subsubsection{Admin Endpoints}
\begin{table*}[t]
\caption{Administrative endpoints.}
\label{tab:admin}
\begin{tabularx}{\textwidth}{l >{\ttfamily}l X}
\toprule
\textbf{Method} & \textbf{Path} & \textbf{Purpose}\\
\midrule
GET & \texttt{/admin/rbac/\{role\}} & Role permissions summary\\
GET & \texttt{/admin/features} & Feature flags + tuning params\\
POST & \texttt{/admin/features/reload} & Reload features.yaml\\
GET & \texttt{/admin/models} & List available LLM models + active\\
POST & \texttt{/admin/models/switch} & Switch active model (persists to .env)\\
GET & \texttt{/admin/enforcement-policy} & Full enforcement policy\\
GET & \texttt{/admin/escalations} & Escalation queue\\
POST & \texttt{/admin/escalations/\{id\}/resolve} & Approve/reject escalation\\
GET & \texttt{/admin/audit/summary} & Audit log summary\\
GET & \texttt{/admin/chain-cache} & Chain cache stats\\
DELETE & \texttt{/admin/chain-cache} & Clear chain cache\\
GET & \texttt{/engagement/scope} & Current engagement scope config\\
\bottomrule
\end{tabularx}
\end{table*}
% =================== 5. NON-FUNCTIONAL REQUIREMENTS ===================
\section{Non-Functional Requirements}

\subsection{Performance Requirements}

\subsubsection{Latency (Measured on CPU-only Hardware)}
\begin{table*}[t]
\caption{Latency metrics.}
\label{tab:latency}
\begin{tabularx}{\textwidth}{l X X X}
\toprule
\textbf{Metric} & \textbf{Raw LLM} & \textbf{Full Pipeline} & \textbf{Notes}\\
\midrule
p50 & 11.5--13.5 s & 16.5 s & Pipeline adds $\sim$3--5 s validation overhead\\
p95 & 16.5--63.3 s & 70.9 s & Dominated by Tier-2 spikes on uncertain inputs\\
p99 & 17.4--93.7 s & $\sim$150 s & Rare Tier-2 chains\\
Mean & 13.4--25.3 s & 28.8 s & Weighted by Tier-1 majority\\
\bottomrule
\end{tabularx}
\end{table*}

\reqid{Requirement NFR-PERF-001:} Pipeline validation overhead (S5--S7) shall not exceed 5 seconds for single-sample requests.

\reqid{Requirement NFR-PERF-002:} Adversarial detection (M1) shall complete in $<$ 5 ms.

\reqid{Requirement NFR-PERF-003:} The semantic cache lookup shall complete in $<$ 50 ms for up to 512 entries.

\textbf{Note:} All latency figures reflect \textbf{CPU-only inference} on an ASUS VivoBook X513EA (Intel i5-1135G7, 4C/8T, 24 GB DDR4). GPU deployment would significantly reduce Tier-1 and Tier-2 latencies.

\subsubsection{Throughput}
\begin{table}[h]
\caption{Throughput limits.}
\label{tab:throughput}
\begin{tabularx}{\columnwidth}{l X}
\toprule
\textbf{Metric} & \textbf{Value}\\
\midrule
Rate limit & 30 requests/minute per IP\\
Concurrent sessions & 1024\\
Cache capacity & 256 (inference) + 512 (semantic)\\
\bottomrule
\end{tabularx}
\end{table}

\subsection{Safety Requirements}
\reqid{Requirement NFR-SAF-001:} The pipeline shall be fail-closed: unknown roles are denied by absence from permit lists; schema violations hard-fail.

\reqid{Requirement NFR-SAF-002:} The pipeline bias shall be toward the safe cell: every observed false positive is an over-block (false rejection), never a dangerous misclassification.

\reqid{Requirement NFR-SAF-003:} Exploitation and password attack intents shall require at least the \texttt{operator} role.

\reqid{Requirement NFR-SAF-004:} Shell metacharacters in target/modifiers shall be blocked regardless of RBAC role.

\reqid{Requirement NFR-SAF-005:} The pipeline shall produce ESCALATE findings (not ALLOW) when enforcement severity is \texttt{"escalate"}, deferring to human review.

\subsection{Security Requirements}

\subsubsection{Input Security}
\reqid{Requirement NFR-SEC-001:} The pipeline shall defend against 4 adversary capabilities:
\begin{enumerate}
\item \textbf{T1 --- Prompt injection:} Lexical injection patterns, structural anomaly detection, LLM self-rejection.
\item \textbf{T2 --- Privilege abuse:} Pre-inference role fast-fail, post-inference role$\times$intent policy.
\item \textbf{T3 --- Model hallucination:} Schema/regex/network/scope validation stack.
\item \textbf{T4 --- Out-of-scope targets:} CIDR containment, RFC-1918 check, scope mode enforcement.
\end{enumerate}

\subsubsection{Data Security}
\reqid{Requirement NFR-SEC-002:} Raw user input shall NEVER be stored in plaintext. Only SHA-256 hashes shall be recorded in the audit log.

\reqid{Requirement NFR-SEC-003:} API authentication shall use the \texttt{X-API-Key} header. The default development key (\texttt{dev\_insecure\_key}) shall NOT be used in production.

\reqid{Requirement NFR-SEC-004:} CORS shall be restricted to specific origins in production (current default: allow all, development mode only).

\subsubsection{Injection Defense Layers}
\begin{table}[h]
\caption{Injection defense layers.}
\label{tab:defense}
\begin{tabularx}{\columnwidth}{l >{\raggedright\arraybackslash}p{0.28\columnwidth} X}
\toprule
\textbf{Layer} & \textbf{Target} & \textbf{Mechanism}\\
\midrule
M1 AdversarialDetector & Raw input & 15 lexical + 7 structural + 3 heuristic patterns\\
S7 ScopeGuard & Raw input & 7 injection regexes (independent of M1)\\
S4 JSONParser & LLM output & Strip thought tags, fences\\
S5 SchemaValidator & Parsed JSON & Pydantic \texttt{extra="forbid"} (no extra fields)\\
S6 RegexValidator & Validated schema & Shell metacharacter scan\\
\bottomrule
\end{tabularx}
\end{table}

\subsection{Reliability and Availability}
\reqid{Requirement NFR-REL-001:} Non-critical middleware failures (session context, adversarial detection, cache store) shall degrade gracefully without aborting the pipeline.

\reqid{Requirement NFR-REL-002:} The audit logger shall be append-only with no delete operations.

\reqid{Requirement NFR-REL-003:} The semantic cache shall self-heal via grounding-hash invalidation when reference datasets change.

\reqid{Requirement NFR-REL-004:} Docker health checks shall verify IRE (30s interval) and Ollama (15s interval) connectivity.

\subsection{Maintainability}
\reqid{Requirement NFR-MNT-001:} All pipeline stages and middleware shall be individually configurable via feature flags in \texttt{config/features.yaml}.

\reqid{Requirement NFR-MNT-002:} Enforcement policy shall be environment-configurable (research/strict/production) without code changes.

\reqid{Requirement NFR-MNT-003:} The test suite shall achieve $\geq$ 80\% line coverage (current: 434 unit + 48 integration = 482 tests).

\subsection{Portability}
\reqid{Requirement NFR-PORT-001:} The IRE shall run on Linux, macOS, and Windows with Python 3.11+.

\reqid{Requirement NFR-PORT-002:} Docker deployment shall use the \texttt{python:3.11-slim} base image.

\reqid{Requirement NFR-PORT-003:} All configuration shall be environment-variable-driven via the \texttt{.env} file or Docker environment.
% =================== 6. DATA REQUIREMENTS ===================
\section{Data Requirements}

\subsection{Logical Data Model}

\subsubsection{Core Entities}
\begin{lstlisting}
IntentSchema
+-- intent: IntentType (9 values)
+-- target: Target {type: TargetType, value: str}
+-- ports: List[int]
+-- modifiers: List[str]
+-- cve_ids: List[str]
+-- tool_hint: Optional[str]
+-- schedule: Optional[str]
+-- confidence: float [0.0, 1.0]
+-- rejection_reason: Optional[str]

PlannerContract (extends IntentSchema)
+-- sub_intent: Optional[SubIntentType]
+-- scope_warnings: List[str]
+-- executability predicates

IREResponseV2 (extends IntentSchema)
+-- validation_findings: List[ValidationFinding]
+-- uncertainty metadata
+-- session context
+-- XAI metadata (optional)
\end{lstlisting}

\subsubsection{ValidationFinding}
\begin{lstlisting}
ValidationFinding
+-- validator: str
+-- passed: bool
+-- hallucination_class: Optional[HallucinationClass]
+-- detail: str
+-- severity: "block" | "warn" | "escalate"
\end{lstlisting}

\subsubsection{CacheEntry (Semantic Cache)}
\begin{lstlisting}
CacheEntry
+-- key_hash: str (SHA-256)
+-- embedding: List[float] (384-dim)
+-- response: IREResponse
+-- grounding_hash: str (SHA-256)
+-- hit_count: int
+-- created_at: float
+-- last_hit_at: float
\end{lstlisting}

\subsubsection{AuditRecord (SQLite)}
\begin{lstlisting}
audit_log
+-- id: INTEGER PRIMARY KEY
+-- timestamp: TEXT (ISO 8601)
+-- session_id: TEXT (nullable)
+-- raw_input_hash: TEXT (SHA-256)
+-- intent_summary: TEXT
+-- hallucination_class: TEXT
+-- severity: TEXT
+-- enforcement_action: TEXT
+-- status: TEXT ("recorded" | "pending" | "approved" | "rejected")
+-- resolved_by: TEXT
+-- resolved_at: TEXT
+-- resolution_note: TEXT
+-- cached_contract_json: TEXT
\end{lstlisting}

\subsection{Data Retention and Privacy}
\begin{table*}[t]
\caption{Data types, retention, and privacy.}
\label{tab:retention}
\begin{tabularx}{\textwidth}{l X X}
\toprule
\textbf{Data Type} & \textbf{Retention} & \textbf{Privacy}\\
\midrule
Raw user input & Transient (in-memory only) & Never stored\\
Audit log & Append-only, persistent & SHA-256 hashed input only\\
Semantic cache & LRU eviction (512 max) & In-memory + optional disk\\
Inference cache & LRU eviction (256 max) & Disk: \texttt{eval/.ollama\_inference\_cache.json}\\
Session context & 30-minute TTL, LRU eviction & In-memory only, never persisted\\
\bottomrule
\end{tabularx}
\end{table*}

\subsection{Grounding Data}

\subsubsection{CVE Grounding Dataset (\texttt{data/cve\_grounding.json})}
\begin{itemize}
\item \textbf{Purpose:} Reference dataset for CVE validation (Check 2 of RegexValidator);
\item \textbf{Format:} JSON with 20 curated entries plus a \texttt{\_meta} block;
\item \textbf{Schema per entry:} \texttt{\{id, description, affected\_service, severity\}};
\item \textbf{Refresh:} Manual via \texttt{scripts/update\_cve\allowbreak{}grounding.py} (CISA KEV catalog);
\item \textbf{Grounding Hash:} SHA-256 digest computed on each cache lookup; stale entries purged.
\end{itemize}

\subsubsection{Alias Map (\texttt{data/alias\_map.json})}
\begin{itemize}
\item \textbf{Purpose:} Map exploit names and slang to canonical CVE identifiers;
\item \textbf{Entries:} 10 exploit-name to CVE mappings + 8 slang/term normalizations;
\item \textbf{Format:} JSON object \texttt{\{alias: canonical, ...\}};
\item \textbf{Processing:} Word-boundary regex compilation, longest-first matching.
\end{itemize}

\subsection{Configuration Data}
\begin{table*}[t]
\caption{Configuration files.}
\label{tab:config}
\begin{tabularx}{\textwidth}{>{\ttfamily}l l X}
\toprule
\textbf{File} & \textbf{Format} & \textbf{Purpose}\\
\midrule
\texttt{config/settings.py} & pydantic-settings & Environment-driven configuration (30+ fields)\\
\texttt{config/features.yaml} & YAML & Feature flags (10) + tuning parameters (11)\\
\texttt{config/rbac\_policy.py} & Python & Role to intent permission table\\
\texttt{config/enforcement\_policy.py} & Python & HallucinationClass to action mapping\\
\texttt{config/lora\_config.yaml} & YAML & Phase-2 fine-tuning plan\\
\bottomrule
\end{tabularx}
\end{table*}

% =================== 7. VERIFICATION AND VALIDATION ===================
\section{Verification and Validation}

\subsection{Test Suite}
\begin{table*}[t]
\caption{Test suite summary.}
\label{tab:testsuite}
\begin{tabularx}{\textwidth}{l l p{3cm} X}
\toprule
\textbf{Category} & \textbf{Count} & \textbf{Files} & \textbf{Notable Coverage}\\
\midrule
Unit tests & 434 & 22 files under \texttt{tests/unit/} & adversarial (42), sub\_intent (46), regex (40), semantic\_cache (37), session\_context (35)\\
Integration tests & 48 & 5 files under \texttt{tests/integration/} & API v1 (12), API v2 (33), audit escalation (3)\\
\textbf{Total} & \textbf{482} & \textbf{27 files} & \\
\bottomrule
\end{tabularx}
\end{table*}

Key behavioral assertions:
\begin{itemize}
\item Tier-2 gating calls the LLM exactly once when confident (\texttt{test\_self\_consistency\_entropy.py});
\item Near-duplicate cache isolation and grounding invalidation (\texttt{test\_semantic\_cache\_adversarial.py});
\item IPv6 and production-mode edge coverage (\texttt{test\_uncovered\_source\_behaviors.py}).
\end{itemize}

\subsection{Evaluation Harness}
\begin{table*}[t]
\caption{Evaluation harness scripts.}
\label{tab:evalharness}
\begin{tabularx}{\textwidth}{>{\ttfamily}l X X}
\toprule
\textbf{Script} & \textbf{Purpose} & \textbf{Path}\\
\midrule
\texttt{diagnose.py} & Full-pipeline runner $\rightarrow$ JSONL debug rows + per-category report & \texttt{eval/diagnose.py}\\
\texttt{metrics.py} & Decision/accuracy/catch/cache metrics & \texttt{eval/metrics.py}\\
\texttt{category\_metrics.py} & NormalizedRow-based metrics engine (canonical) & \texttt{eval/category\_metrics.py}\\
\texttt{run\_baseline.py} & Raw unconstrained LLM baseline scorer & \texttt{eval/run\_baseline.py}\\
\texttt{exp1\_raw\_vs\_pipeline.py} & Exp1: raw vs. pipeline, 4 configs & \texttt{eval/Data Collection/}\\
\texttt{exp2\_model\_head\_to\_head.py} & Exp2: Gemma vs. Qwen & \texttt{eval/Data Collection/}\\
\texttt{exp3\_ablation\_study.py} & Exp3: 8-config $\times$ 900-record ablation & \texttt{eval/Data Collection/}\\
\texttt{exp4\_hallucination\_analysis.py} & Exp4: per-class deep-dive (250 records) & \texttt{eval/Data Collection/}\\
\bottomrule
\end{tabularx}
\end{table*}

\subsection{Empirical Results Summary (900-Record Canonical Runs)}

\subsubsection{Experiment 1 --- Raw LLM vs. Full Pipeline}
\begin{table}[h]
\caption{Experiment 1 results.}
\label{tab:exp1}
\begin{tabularx}{\columnwidth}{l X X X X}
\toprule
\textbf{Metric} & \textbf{Raw Gemma} & \textbf{Raw Qwen} & \textbf{Pipeline Gemma} & \textbf{Pipeline Qwen}\\
\midrule
Intent Accuracy & 78.0\% & 69.0\% & \textbf{86.0\%} & \textbf{82.0\%}\\
Hallucination Catch & 0.0\% & 0.0\% & \textbf{72.0\%} & \textbf{64.0\%}\\
False Positive Rate & 32.0\% & 31.0\% & \textbf{14.0\%} & \textbf{15.0\%}\\
\bottomrule
\end{tabularx}
\end{table}

\noindent\textbf{Pipeline lift:} +8pp/+13pp accuracy, +72pp/+64pp catch, $-18$pp/$-16$pp FP.

\subsubsection{Experiment 2 --- Model Head-to-Head}
\begin{table}[h]
\caption{Experiment 2 results.}
\label{tab:exp2}
\begin{tabularx}{\columnwidth}{l X X X}
\toprule
\textbf{Metric} & \textbf{Gemma 4} & \textbf{Qwen 2.5 Coder} & \textbf{Delta}\\
\midrule
Intent Accuracy & 86.0\% & 82.0\% & $-4.0$pp\\
Hallucination Catch & 72.0\% & 64.0\% & $-8.0$pp\\
FP Rate & 14.0\% & 15.0\% & $+1.0$pp\\
p50 Latency & 27,650 ms & 32,755 ms & $+18\%$\\
\bottomrule
\end{tabularx}
\end{table}

\subsubsection{Experiment 3 --- Component Ablation}
\begin{table*}[t]
\caption{Experiment 3 ablation results.}
\label{tab:exp3}
\begin{tabularx}{\textwidth}{l X X X X X}
\toprule
\textbf{Config} & \textbf{DAcc} & \textbf{Catch} & \textbf{FP} & \textbf{CFAcc} & \textbf{Mean Latency}\\
\midrule
\textbf{Full Pipeline} & \textbf{73.5\%} & \textbf{84\%} & 34\% & 87.5\% & 28,805 ms\\
No Scope Guard & 64.0\% & 73\% & 44\% & 83.9\% & 31,404 ms\\
No RBAC Guard & 74.5\% & 64\% & 22\% & 88.2\% & 28,260 ms\\
No Network Validator & 68.0\% & 82\% & 44\% & 83.9\% & 27,450 ms\\
No Hallucination Validator & 65.4\% & 62\% & 44\% & 82.1\% & 26,722 ms\\
Validation Stack Disabled & 64.0\% & 41\% & 13\% & 87.4\% & 25,198 ms\\
Self-Consistency Disabled & 67.5\% & 80\% & 45\% & 85.5\% & 11,809 ms\\
Semantic Cache Disabled & 66.0\% & 79\% & 47\% & 86.8\% & 22,949 ms\\
\bottomrule
\end{tabularx}
\end{table*}

\textbf{Component contribution ranking (by accuracy impact):}
\begin{enumerate}
\item ScopeGuard ($+9.5$pp);
\item Entire Validation Stack ($+9.5$pp);
\item SchemaValidator + RegexValidator ($+8.1$pp);
\item Semantic Cache ($+7.5$pp);
\item Tier-2 Self-Consistency ($+6.0$pp);
\item NetworkArchitectureValidator ($+5.5$pp);
\item RBACGuard ($-1.0$pp).
\end{enumerate}

\subsubsection{Experiment 4 --- Hallucination Deep-Dive (250 Records)}
\begin{table}[h]
\caption{Experiment 4 per-category catch rates.}
\label{tab:exp4}
\begin{tabularx}{\columnwidth}{l X X >{\raggedright\arraybackslash}p{0.15\columnwidth}}
\toprule
\textbf{Category} & \textbf{Total} & \textbf{Caught} & \textbf{Catch \%}\\
\midrule
rbac\_violation & 12 & 12 & \textbf{100.0\%}\\
adversarial\_injection & 45 & 43 & \textbf{95.6\%}\\
fabricated\_parameter & 32 & 28 & 87.5\%\\
target\_type\_mismatch & 32 & 26 & 81.2\%\\
fabricated\_cve & 32 & 25 & 78.1\%\\
contradictory\_action\_target & 32 & 19 & 59.4\%\\
out\_of\_scope & 45 & 21 & \textbf{46.7\%}\\
ambiguous & 20 & 8 & 40.0\%\\
\textbf{Total} & \textbf{250} & \textbf{182} & \textbf{72.8\%}\\
\bottomrule
\end{tabularx}
\end{table}

\textbf{Largest structural gap:} \texttt{llm\_no\_scope\_detection} (35\% of all misses) --- public/out-of-scope targets reach the LLM before any scope gate.

\subsection{Metric Definitions}
\begin{table}[h]
\caption{Metric definitions.}
\label{tab:metrics}
\begin{tabularx}{\columnwidth}{>{\bfseries}l X}
\toprule
\textbf{Metric} & \textbf{Formula}\\
\midrule
DAcc (Decision Accuracy) & \texttt{sum}/\texttt{N\_total} of correct decisions: well\_formed$\rightarrow$ALLOWED+correct fields, flawed$\rightarrow$BLOCKED/ESCALATED\\
Catch (Flawed-Input Catch) & \texttt{|flawed records with BLOCKED/ESCALATED|}/\texttt{N\_flawed}\\
FP (False Positive Rate) & \texttt{|well\_formed records with non-ALLOWED decision|}/\texttt{N\_well\_formed}\\
CFAcc (Contract-Field Accuracy) & \texttt{|ALLOWED records with correct intent+target|}/\texttt{|ALLOWED well\_formed|}\\
\bottomrule
\end{tabularx}
\end{table}
% =================== 8. DEPLOYMENT ===================
\section{Deployment}

\subsection{Docker Deployment (Recommended)}
\begin{lstlisting}
# 1. Clone and configure
cd apps/services/neuroshell-ire
cp .env.example .env
# Edit .env: set IRE_API_KEY, OLLAMA_MODEL

# 2. Build and start
docker-compose up -d --build

# 3. Verify
curl http://localhost:8001/health
\end{lstlisting}

Services:
\begin{table*}[t]
\caption{Docker services.}
\label{tab:docker}
\begin{tabularx}{\textwidth}{>{\ttfamily}l >{\ttfamily}l >{\ttfamily}l >{\raggedright\arraybackslash}X}
\toprule
\textbf{Service} & \textbf{Container} & \textbf{Port} & \textbf{Health Check}\\
\midrule
ire-service & neuroshell-ire & 8001 & \texttt{curl -f http://localhost:8001/health}\\
ollama & neuroshell-ollama & 11434 & \texttt{curl -f http://localhost:11434/api/tags}\\
\bottomrule
\end{tabularx}
\end{table*}

\noindent\textbf{Volumes:} \texttt{ollama\_models} mounted at \texttt{/root/.ollama} (persists downloaded models).

\subsection{Manual Deployment}
\begin{lstlisting}
# 1. Prerequisites
python --version  # 3.11+
ollama --version  # 0.32.x

# 2. Virtual environment
python -m venv .venv
.venv\Scripts\activate  # Windows
# source .venv/bin/activate  # Linux/macOS

# 3. Dependencies
pip install -r requirements.txt
python -m spacy download en_core_web_sm

# 4. Configuration
cp .env.example .env
# Edit .env: set IRE_API_KEY, OLLAMA_MODEL, OLLAMA_BASE_URL

# 5. Pull LLM model
ollama pull gemma4:e4b

# 6. Start server
python scripts/run_dev.py
\end{lstlisting}

\subsection{Configuration Management}

\subsubsection{Environment Variables (.env)}
\begin{table*}[t]
\caption{Environment variables.}
\label{tab:envvars}
\begin{tabularx}{\textwidth}{>{\ttfamily}l >{\ttfamily}l X}
\toprule
\textbf{Variable} & \textbf{Default} & \textbf{Description}\\
\midrule
\texttt{IRE\_API\_KEY} & \texttt{dev\_insecure\_key} & API authentication key\\
\texttt{OLLAMA\_BASE\_URL} & \texttt{http://localhost:11434} & Ollama daemon URL\\
\texttt{OLLAMA\_MODEL} & \texttt{gemma4:e4b} & Active LLM model\\
\texttt{SCOPE\_MODE} & \texttt{research} & Scope enforcement mode\\
\texttt{ENGAGEMENT\_SCOPE} & \texttt{192.168.0.0/16} & Authorised target CIDR\\
\texttt{LOG\_LEVEL} & \texttt{INFO} & Logging level\\
\texttt{PORT} & \texttt{8001} & Server port\\
\bottomrule
\end{tabularx}
\end{table*}

\subsubsection{Feature Flags (config/features.yaml)}
\begin{table}[h]
\caption{Feature flags.}
\label{tab:flags}
\begin{tabularx}{\columnwidth}{>{\ttfamily}l >{\ttfamily}l X}
\toprule
\textbf{Flag} & \textbf{Default} & \textbf{Description}\\
\midrule
\texttt{session\_context} & true & Multi-turn session memory\\
\texttt{rbac} & true & Role-based access control\\
\texttt{adversarial\_detection} & true & Adversarial input detection\\
\texttt{semantic\_cache} & true & Embedding-based response cache\\
\texttt{sub\_intent} & true & Hierarchical sub-intent classification\\
\texttt{schema\_v2} & true & V2 schema routing\\
\texttt{xai} & false & XAI token attribution\\
\texttt{uncertainty\_estimation} & false & Monte Carlo dropout\\
\texttt{continual\_learning} & false & EWC active learning\\
\texttt{speculative\_decoding} & false & Draft model acceleration\\
\bottomrule
\end{tabularx}
\end{table}

% =================== 9. APPENDICES ===================
\appendices

\section{IntentType Enumeration (9 Values)}
\begin{table*}[t]
\caption{Intent types.}
\label{tab:appA}
\begin{tabularx}{\textwidth}{>{\ttfamily}l X >{\raggedright\arraybackslash}p{0.4\textwidth}}
\toprule
\textbf{Value} & \textbf{Category} & \textbf{Description}\\
\midrule
NETWORK\_SCAN & Active Recon & Port scanning, host discovery, OS detection\\
VULNERABILITY\_AUDIT & Assessment & CVE-specific or service-specific vulnerability checks\\
DIRECTORY\_BRUTEFORCE & Enumeration & Web directory/file brute-forcing\\
SERVICE\_ENUMERATION & Enumeration & Service version detection, banner grabbing\\
EXPLOITATION & Attack & Remote code execution, privilege escalation\\
PASSWORD\_ATTACK & Attack & Brute force, credential stuffing, password spray\\
PASSIVE\_RECON & Recon & DNS lookup, WHOIS, OSINT gathering\\
AMBIGUOUS & Special & Confident $<$ 0.5; requires clarification\\
REJECTED & Special & Command rejected (injection, policy, etc.)\\
\bottomrule
\end{tabularx}
\end{table*}

\section{TargetType Enumeration (6 Values)}
\begin{table*}[t]
\caption{Target types.}
\label{tab:appB}
\begin{tabularx}{\textwidth}{>{\ttfamily}l X X}
\toprule
\textbf{Value} & \textbf{Format} & \textbf{Example}\\
\midrule
IP & IPv4 address & \texttt{192.168.1.1}\\
SUBNET & CIDR notation & \texttt{192.168.1.0/24}\\
DOMAIN & Domain name & \texttt{example.com}\\
URL & HTTP/HTTPS URL & \texttt{https://example.com/api}\\
HOSTNAME & Single-label hostname & \texttt{webserver01}\\
UNKNOWN & Empty/unspecified & (empty string)\\
\bottomrule
\end{tabularx}
\end{table*}

\section{SubIntentType Enumeration (27 Values)}
\begin{table}[h]
\caption{Sub-intent types by parent intent.}
\label{tab:appC}
\begin{tabularx}{\columnwidth}{>{\ttfamily}l >{\ttfamily}X}
\toprule
\textbf{Parent Intent} & \textbf{Sub-Intents}\\
\midrule
NETWORK\_SCAN & ICMP\_DISCOVERY, SYN\_STEALTH, UDP\_SWEEP, VERSION\_DETECT, OS\_DETECT, FULL\_PORT\\
EXPLOITATION & RCE, PRIV\_ESC, LATERAL, PERSISTENCE\\
VULNERABILITY\_AUDIT & CVE\_SPECIFIC, SERVICE\_SPECIFIC, FULL\_SCAN\\
SERVICE\_ENUMERATION & SMB, SNMP, VERSION\_DETECT, BANNER\_GRAB\\
PASSWORD\_ATTACK & SPRAY, CREDENTIAL\_STUFFING, BRUTE\_FORCE\\
PASSIVE\_RECON & DNS, WHOIS, OSINT\\
DIRECTORY\_BRUTEFORCE & API, FILES, WEB\\
\bottomrule
\end{tabularx}
\end{table}

\section{HallucinationClass Enumeration (6 Values)}
\begin{table}[h]
\caption{Hallucination classes and enforcement.}
\label{tab:appD}
\begin{tabularx}{\columnwidth}{>{\ttfamily}l X X}
\toprule
\textbf{Class} & \textbf{Detection} & \textbf{Default Enforcement}\\
\midrule
FABRICATED\_CVE & S5 + S6 (syntax) & BLOCK\\
TARGET\_TYPE\_MISMATCH & S5 + S6 (format) + S6B & BLOCK\\
CONTRADICTORY\_ACTION\_TARGET & S5 (schema) & BLOCK\\
FABRICATED\_PARAMETER & S5 + S6 (port/metachar) & BLOCK\\
UNGROUNDED\_CONFIDENCE & S5 (confidence) + S6 (grounding) & WARN\\
OUT\_OF\_SCOPE\_TARGET & S6B (scope) + S7 (RFC-1918) & Dynamic (BLOCK/WARN)\\
\bottomrule
\end{tabularx}
\end{table}

\section{RBAC Policy Table}
\begin{table*}[t]
\caption{RBAC policy matrix.}
\label{tab:appE}
\begin{tabularx}{\textwidth}{>{\ttfamily}l c c c c}
\toprule
\textbf{Intent} & \textbf{viewer} & \textbf{analyst} & \textbf{operator} & \textbf{admin}\\
\midrule
NETWORK\_SCAN & --- & \checkmark & \checkmark & \checkmark\\
VULNERABILITY\_AUDIT & --- & \checkmark & \checkmark & \checkmark\\
DIRECTORY\_BRUTEFORCE & --- & --- & \checkmark & \checkmark\\
SERVICE\_ENUMERATION & --- & \checkmark & \checkmark & \checkmark\\
EXPLOITATION & --- & --- & \checkmark & \checkmark\\
PASSWORD\_ATTACK & --- & --- & \checkmark & \checkmark\\
PASSIVE\_RECON & --- & \checkmark & \checkmark & \checkmark\\
AMBIGUOUS & --- & \checkmark & \checkmark & \checkmark\\
REJECTED & --- & \checkmark & \checkmark & \checkmark\\
\bottomrule
\end{tabularx}
\end{table*}

\noindent\textbf{Sensitive intents (require operator+):} EXPLOITATION, PASSWORD\_ATTACK, DIRECTORY\_BRUTEFORCE.
\section{Enforcement Policy Table}
\begin{table*}[t]
\caption{Enforcement policy by environment mode.}
\label{tab:appF}
\begin{tabularx}{\textwidth}{>{\ttfamily}l l l l}
\toprule
\textbf{HallucinationClass} & \textbf{research} & \textbf{strict} & \textbf{production}\\
\midrule
FABRICATED\_PARAMETER & BLOCK & BLOCK & BLOCK\\
CONTRADICTORY\_ACTION\_TARGET & BLOCK & BLOCK & BLOCK\\
TARGET\_TYPE\_MISMATCH & BLOCK & BLOCK & BLOCK\\
FABRICATED\_CVE & BLOCK & BLOCK & BLOCK\\
UNGROUNDED\_CONFIDENCE & WARN & WARN & WARN\\
OUT\_OF\_SCOPE\_TARGET & WARN & BLOCK & BLOCK\\
\bottomrule
\end{tabularx}
\end{table*}

\section{Adversarial Detection Patterns}

\subsection{Layer A: Lexical (15 patterns)}
\begin{table*}[t]
\caption{Layer A lexical patterns.}
\label{tab:appG1}
\begin{tabularx}{\textwidth}{>{\ttfamily}X >{\raggedright\arraybackslash}p{0.08\textwidth}}
\toprule
\textbf{Pattern} & \textbf{Weight}\\
\midrule
\texttt{ignore (all~)?(previous|prior|above) instructions?} & 0.9\\
\texttt{forget (everything|all|your instructions?)} & 0.9\\
\texttt{you are now (?!scanning|checking|running)} & 0.8\\
\texttt{act as (a~)?(different|new|another|evil|unrestricted)} & 0.8\\
\texttt{jailbreak} & 1.0\\
\texttt{dan mode} & 1.0\\
\texttt{prompt injection} & 0.9\\
\texttt{(system|assistant|user):\textbackslash{}s*you (must|shall|will)} & 0.85\\
\texttt{$<$|im\_(start|end)|$>$} & 1.0\\
\texttt{$<$|think|$>$} & 0.95\\
\texttt{[INST]} or \texttt{[/INST]} & 0.9\\
\texttt{\#\#\# (instruction|system|human|assistant)} & 0.85\\
\texttt{disregard (all~)?(prior|previous|your)} & 0.85\\
\texttt{new instructions:} & 0.8\\
\texttt{override (previous~)?(instructions?|rules?|guidelines?)} & 0.85\\
\bottomrule
\end{tabularx}
\end{table*}

\subsection{Layer B: Structural (7 patterns + 3 heuristics)}
\begin{table*}[t]
\caption{Layer B structural patterns.}
\label{tab:appG2}
\begin{tabularx}{\textwidth}{>{\ttfamily}X >{\raggedright\arraybackslash}p{0.08\textwidth}}
\toprule
\textbf{Pattern} & \textbf{Weight}\\
\midrule
\texttt{(\textbackslash{}x[0-9a-f]\{2\})\{4,\}} & 0.6\\
\texttt{(\%[0-9a-f]\{2\})\{4,\}} & 0.5\\
\texttt{base64:\textbackslash{}s*[A-Za-z0-9+/]\{20,\}} & 0.7\\
\texttt{$<$script[\textbackslash{}s$>$]} & 0.8\\
\texttt{\textbackslash{}\$\{[\^{}]\}\{0,50\}\textbackslash{}\}} & 0.6\\
\texttt{eval( |exec(|\_\_import\_\_()} & 0.9\\
\texttt{rm -rf|del /[sq]|format c:} & 0.8\\
\bottomrule
\end{tabularx}
\end{table*}

\begin{table}[h]
\caption{Layer B heuristics.}
\label{tab:appG3}
\begin{tabularx}{\columnwidth}{l X X}
\toprule
\textbf{Heuristic} & \textbf{Threshold} & \textbf{Weight}\\
\midrule
Special-char density & $>$ 30\% (for inputs $>$ 20 chars) & 0.4\\
Single token length & $>$ 100 chars & 0.3\\
Repeated n-gram & $>$ 3 occurrences of 3/4/5-grams ($\geq$ 10 chars) & 0.35\\
\bottomrule
\end{tabularx}
\end{table}

\section{Configuration Parameters (Complete)}
\begin{table*}[t]
\caption{All configuration parameters.}
\label{tab:appH}
\begin{tabularx}{\textwidth}{>{\ttfamily}l X l l}
\toprule
\textbf{Parameter} & \textbf{Default} & \textbf{Type} & \textbf{Source}\\
\midrule
\texttt{ire\_api\_key} & \texttt{dev\_insecure\_key} & str & .env\\
\texttt{port} & 8001 & int & .env\\
\texttt{ollama\_base\_url} & \texttt{http://localhost:11434} & str & .env\\
\texttt{ollama\_model} & \texttt{gemma4:e4b} & str & .env\\
\texttt{ollama\_timeout\_seconds} & 60 & int & .env\\
\texttt{inference\_temperature} & 0.1 & float & .env\\
\texttt{inference\_max\_tokens} & 2048 & int & .env\\
\texttt{self\_consistency\_trigger\_threshold} & 0.7 & float & .env\\
\texttt{self\_consistency\_samples} & 5 & int & .env\\
\texttt{self\_consistency\_temperature} & 0.7 & float & .env\\
\texttt{entropy\_low\_threshold} & 0.3 & float & .env\\
\texttt{entropy\_high\_threshold} & 0.8 & float & .env\\
\texttt{scope\_mode} & research & literal & .env\\
\texttt{engagement\_scope} & \texttt{192.168.0.0/16} & str & .env\\
\texttt{engagement\_name} & Lab Environment & str & .env\\
\texttt{engagement\_mode} & warn & literal & .env\\
\texttt{max\_cidr\_prefix} & 30 & int & .env\\
\texttt{min\_cidr\_prefix} & 32 & int & .env\\
\texttt{lru\_cache\_max\_size} & 256 & int & .env\\
\texttt{planner\_enabled} & false & bool & .env\\
\texttt{planner\_url} & \texttt{http://127.0.0.1:8002} & str & .env\\
\texttt{planner\_api\_key} & \texttt{neuroshell-c2-secret-key} & str & .env\\
\texttt{planner\_timeout\_seconds} & 300 & int & .env\\
\texttt{component3\_enabled} & false & bool & .env\\
\texttt{component3\_url} & \texttt{http://127.0.0.1:8003} & str & .env\\
\texttt{component4\_enabled} & false & bool & .env\\
\texttt{component4\_url} & \texttt{http://127.0.0.1:8004} & str & .env\\
\texttt{session\_ttl\_minutes} & 30 & int & features.yaml\\
\texttt{session\_max\_turns} & 8 & int & features.yaml\\
\texttt{session\_max\_sessions} & 1024 & int & features.yaml\\
\texttt{semantic\_cache\_threshold} & 0.82 & float & features.yaml\\
\texttt{semantic\_cache\_max\_size} & 512 & int & features.yaml\\
\texttt{adversarial\_classifier\_threshold} & 0.7 & float & features.yaml\\
\texttt{sub\_intent\_confidence\_threshold} & 0.6 & float & features.yaml\\
\bottomrule
\end{tabularx}
\end{table*}

\section{Exception Hierarchy}
\begin{lstlisting}
IREError (base)
+-- JSONParseError          (stage: "json_parser")
+-- SchemaValidationError   (stage: "schema_validator")
+-- RegexValidationError    (stage: "regex_validator")
+-- NetworkValidationError  (stage: "network_validator")
|   +-- extra: validation_type
+-- ScopeError              (stage: "scope_guard")
+-- InferenceError          (stage: "inference")
+-- RBACError               (stage: "rbac_guard")
    +-- extra: role, denied_intent
\end{lstlisting}

\section{Dependency Versions}
\begin{table*}[t]
\caption{Package dependency versions.}
\label{tab:appJ}
\begin{tabularx}{\textwidth}{>{\ttfamily}l >{\ttfamily}l X}
\toprule
\textbf{Package} & \textbf{Version} & \textbf{Purpose}\\
\midrule
\texttt{fastapi} & 0.115.0 & Async REST framework\\
\texttt{uvicorn[standard]} & 0.30.6 & ASGI server\\
\texttt{pydantic} & $\geq$2.9, $<$3.0 & Schema validation\\
\texttt{pydantic-settings} & 2.4.0 & Environment config\\
\texttt{jsonschema} & 4.23.0 & JSON Schema support\\
\texttt{spacy} & 3.7.6 & NLP preprocessing\\
\texttt{ollama} & 0.3.3 & LLM client\\
\texttt{httpx} & 0.27.2 & HTTP client\\
\texttt{structlog} & 24.4.0 & Structured logging\\
\texttt{sentence-transformers} & 3.0.1 & Semantic cache embeddings\\
\texttt{slowapi} & 0.1.9 & Rate limiting\\
\texttt{pytest} & 8.3.2 & Test framework\\
\texttt{pytest-asyncio} & 0.23.8 & Async test support\\
\bottomrule
\end{tabularx}
\end{table*}

\section{Hashing Algorithms}
\begin{table*}[t]
\caption{Hash algorithm uses.}
\label{tab:appK}
\begin{tabularx}{\textwidth}{l l X}
\toprule
\textbf{Use} & \textbf{Algorithm} & \textbf{Input}\\
\midrule
Audit input privacy & SHA-256 & Raw user input\\
Semantic cache exact key & SHA-256 & Lowercased, stripped input text\\
Inference cache key & SHA-256 & \texttt{model\_name|enriched\_input}\\
Grounding invalidation & SHA-256 & Concatenation of \texttt{cve\_grounding.json} + \texttt{alias\_map.json}\\
\bottomrule
\end{tabularx}
\end{table*}

\vspace{1em}
\noindent\textit{--- End of Document ---}

\end{document}
